% !TeX root = ../papers/nc-p3.tex
\section{Bell曲面の全流入窓の周期二部分の四成分}
\label{sec:split-surface-periods}
\begin{theorem}\label{thm:split-surface-periods}
前節の曲面の一般の閉軌道の四補完のいずれでも、閉軌道の最小周期は2である。
閉軌道に至る非自明な全窓のうち、周期2の軌跡は同型を除いて
次元 $4,2,2,2$ の四つの既約成分を持ち、それ以外の点は全て最小周期4である。

具体的に、重み1の四座標を $b$、重み2の三座標を $c$ とし、
$\rho$ を下記の補正二移動、$k$ を $1,Z,X,XZ$ の四つの補正四元群元とする。
$E_\pm(k\rho)$ を該当する座標空間での固有空間と書く。
$b\ne0$ の周期2条件は、ある $k$ で $b\in E_+(k\rho)\cup E_-(k\rho)$ かつ
$c\in E_+(k\rho)$ となることである。$b=0$ なら常に周期2である。
\end{theorem}
\begin{proof}[一般点での二移動]
$P$ の一巡回あるいはその相補逆の一巡回が元のBell四数と局所同値でなければ、
軸交換も許して一移動同型は生じない。$(u,v)=(4,7)$ で
$(1,-28,4,7)$ に対する比較相手は $(-15,-12,20,-9)$ と $(15,12,20,-9)/28$ となり、
どちらも射影的な偶符号変更・置換で元に戻らない。
ゆえに一般点で一移動はなく、閉極限の一意性からその上の全窓でもない。

$\kappa=-(u^2-1)(v^2-1)/(4uv)$ とする。二移動で四テンソルの
六ブロックのうち法線位置が $01$ のものは $\kappa$ 倍、$23$ は $\kappa^{-1}$ 倍、
他の四つは一定である。二移動の後で位置別の法線文字だけに
$(1/s,1/s,s,s)$ を掛けると $w_2$ を固定する（$s^2=\kappa$）。これを $\rho$ とする。
$Z$ は全二文字のビット符号、$X$ は全ビット反転の後に法線だけに
$(\iota,\iota,-\iota,-\iota)$ を掛ける補正とし、$\iota^2=-1$ とする。
これらも $w_2$ を固定する。テンソル全体で $X,Z$ は可換、
$X^2$ は法線数 $n$ の空間で $(-1)^n$、$Z^2=\rho^2=1$ である。
$\rho Z=Z\rho$、$\rho X=(-1)^n X\rho$ なので、四つの $k\rho$ は全て対合となる。
この有限群と法線倍率を含めて前節の同変な二つの切断を選べば、
同じ作用が切片の $(b,c)$ に線型に作用する。

重み1側は三法線の一次核をせん断で割った四空間に等しく、
重み2側は四法線の二つの射影の共通核である。
前者の $Z$ 正負固有空間は各二次元で、各 $k\rho$ の正負固有空間も各二次元である。
後者では $k=hZ$（$h$ は補完の添字）のとき $k\rho$ は恒等、
他の三つでは正固有空間は一次元、負固有空間は二次元となる。
\path{computations/verify_split_surface_periods.py} は四つの切断の倍率を一般体で、
固有空間を $u=2,v=3$, 素数109、$s=\iota=33$ で計算する。
ここで一次核・せん断・四法線核の次元は一般値を保つ。階数が一定の近傍では
これらの格子は局所自由である。2を可逆にすると対合の二射影は直和因子となるので
固有空間の次元が一般体にも戻る。$s$ のもう一方の選択は重み1側で符号を交換するだけである。
\end{proof}
\begin{proof}[全座標での判定と成分数]
二移動後に $k$ と法線倍率を用いて戻ればよい。
$b\ne0$ のとき $k\rho$ の $b$ 上の固有値は $1$ または $-1$ であり、
その平方が $c$ 上の倍率となるため $c$ は正固有空間に入る。
$b=0$ は $k=hZ$ で全て含まれる。$X$ は重み1上で $k\rho$ と反交換し、
重み2上では可換なので二符号は同型になる。
各 $k$ の像は従って一つの既約な重み付き線型部分空間の有限商であり、
その次元は $2+\dim E_+(k\rho)|_c-1$ である。

相異なる $k$ の像の一般の $b$ の方向は異なる。
$Z$ の各固有二平面で $X$ の二固有線は $\rho$ によって交換されるので、
$k\rho$ の各固有平面は $Z$ の二平面それぞれから一本の固有線を取る。
$k$ を $kZ$ に替えると一方だけの符号が反転する。
$k$ を $kX$ または $kXZ$ に替えると、各 $Z$ 固有二平面で元の対合と反交換するため
固有線は互いに異なる。四元群での共役は各 $k\rho$ の符号だけを変える。
従って小次元の三つの像も他の像に含まれず、表示した四成分で尽きる。
\end{proof}
